\input{preamble/preamble.tex}

\geometry{margin=0.85in}
\AtBeginDocument{\small}

\newcommand{\ExamNameField}{\noindent\textbf{Name:}\ \rule{0.7\linewidth}{0.4pt}\par\medskip}

\begin{document}
	\ExamTitleBlock{11th grade}{Term 2 Learning Evidence 2: C7 Activity Solution}{}
	
	\section*{1. Local Transport Satisfaction}
		
	A city transport office surveys \textbf{100 commuters} about whether they are satisfied with the reliability of the bus system. In the sample, \textbf{58 commuters} report being satisfied.
	
	Construct and interpret a \textbf{90\% confidence interval} for the population proportion of commuters who are satisfied with the bus system.
	
	\subsubsection*{(C7) Constructs the confidence interval for a proportion}
	To construct the confidence interval for a proportion, we identify the population proportion, compute \(\hat{p}\), verify large-sample conditions, select the correct \(z^*\), compute \(SE\) and the margin of error, construct the interval, and interpret the result in context.
	
	\textbf{Step 1: Identify the population parameter and the sample proportion.}
	Let \(p\) be the population proportion of commuters who are satisfied. The sample proportion is
	\[
	\hat{p} = \frac{x}{n} = \frac{58}{100} = 0.58
	\]
	Large-sample conditions:
	\[
	n\hat{p} = 100(0.58) = 58 \ge 10, \qquad n(1-\hat{p}) = 100(0.42) = 42 \ge 10
	\]
	
	\textbf{Step 2: Compute the standard error of the proportion.}
	\[
	SE = \sqrt{\frac{\hat{p}(1-\hat{p})}{n}} = \sqrt{\frac{0.58(0.42)}{100}} \approx 0.0494
	\]
	
	\textbf{Step 3: Compute the margin of error (90\% confidence, \(z^* = 1.645\)).}
	\[
	E = z^* \cdot SE = 1.645(0.0494) \approx 0.0812
	\]
	
	\textbf{Step 4: Construct the confidence interval.}
	\[
	\hat{p} \pm E = 0.58 \pm 0.0812 \Rightarrow [0.499,\ 0.661]
	\]
	
	\textbf{Step 5: Interpret the confidence interval in context.}
	We are 90\% confident that between about 49.9\% and 66.1\% of all commuters in the city are satisfied with the bus system.
	
	\section*{2. Digital Payment Adoption}
		
	A business association surveys \textbf{180 small retailers} about whether they accept digital payments. In the sample, \textbf{81 retailers} accept digital payments.
	
	Construct and interpret a \textbf{95\% confidence interval} for the population proportion of small retailers who accept digital payments.
	
	\subsubsection*{(C7) Constructs the confidence interval for a proportion}
	To construct the confidence interval for a proportion, we identify the population proportion, compute \(\hat{p}\), verify large-sample conditions, select the correct \(z^*\), compute \(SE\) and the margin of error, construct the interval, and interpret the result in context.
	
	\textbf{Step 1: Identify the population parameter and the sample proportion.}
	Let \(p\) be the population proportion of small retailers who accept digital payments. The sample proportion is
	\[
	\hat{p} = \frac{x}{n} = \frac{81}{180} = 0.45
	\]
	Large-sample conditions:
	\[
	n\hat{p} = 180(0.45) = 81 \ge 10, \qquad n(1-\hat{p}) = 180(0.55) = 99 \ge 10
	\]
	
	\textbf{Step 2: Compute the standard error of the proportion.}
	\[
	SE = \sqrt{\frac{\hat{p}(1-\hat{p})}{n}} = \sqrt{\frac{0.45(0.55)}{180}} \approx 0.0371
	\]
	
	\textbf{Step 3: Compute the margin of error (95\% confidence, \(z^* = 1.96\)).}
	\[
	E = z^* \cdot SE = 1.96(0.0371) \approx 0.0727
	\]
	
	\textbf{Step 4: Construct the confidence interval.}
	\[
	\hat{p} \pm E = 0.45 \pm 0.0727 \Rightarrow [0.377,\ 0.523]
	\]
	
	\textbf{Step 5: Interpret the confidence interval in context.}
	We are 95\% confident that between about 37.7\% and 52.3\% of all small retailers accept digital payments.
	
	\section*{3. Water-Saving Appliance Use}
		
	An environmental agency surveys \textbf{250 households} to estimate the share that use water-saving appliances. In the sample, \textbf{190 households} report using such appliances.
	
	Construct and interpret a \textbf{95\% confidence interval} for the population proportion of households using water-saving appliances.
	
	\subsubsection*{(C7) Constructs the confidence interval for a proportion}
	To construct the confidence interval for a proportion, we identify the population proportion, compute \(\hat{p}\), verify large-sample conditions, select the correct \(z^*\), compute \(SE\) and the margin of error, construct the interval, and interpret the result in context.
	
	\textbf{Step 1: Identify the population parameter and the sample proportion.}
	Let \(p\) be the population proportion of households that use water-saving appliances. The sample proportion is
	\[
	\hat{p} = \frac{x}{n} = \frac{190}{250} = 0.76
	\]
	Large-sample conditions:
	\[
	n\hat{p} = 250(0.76) = 190 \ge 10, \qquad n(1-\hat{p}) = 250(0.24) = 60 \ge 10
	\]
	
	\textbf{Step 2: Compute the standard error of the proportion.}
	\[
	SE = \sqrt{\frac{\hat{p}(1-\hat{p})}{n}} = \sqrt{\frac{0.76(0.24)}{250}} \approx 0.0270
	\]
	
	\textbf{Step 3: Compute the margin of error (95\% confidence, \(z^* = 1.96\)).}
	\[
	E = z^* \cdot SE = 1.96(0.0270) \approx 0.0529
	\]
	
	\textbf{Step 4: Construct the confidence interval.}
	\[
	\hat{p} \pm E = 0.76 \pm 0.0529 \Rightarrow [0.707,\ 0.813]
	\]
	
	\textbf{Step 5: Interpret the confidence interval in context.}
	We are 95\% confident that between about 70.7\% and 81.3\% of households use water-saving appliances.
	
	\section*{4. Export Quality Compliance}
		
	A trade ministry inspects \textbf{320 export shipments} and finds that \textbf{68} fail to meet a quality standard. The ministry wants to estimate the proportion of shipments that fail the standard.
	
	Construct and interpret a \textbf{98\% confidence interval} for the population proportion of export shipments that fail the quality standard.
	
	\subsubsection*{(C7) Constructs the confidence interval for a proportion}
	To construct the confidence interval for a proportion, we identify the population proportion, compute \(\hat{p}\), verify large-sample conditions, select the correct \(z^*\), compute \(SE\) and the margin of error, construct the interval, and interpret the result in context.
	
	\textbf{Step 1: Identify the population parameter and the sample proportion.}
	Let \(p\) be the population proportion of export shipments that fail the quality standard. The sample proportion is
	\[
	\hat{p} = \frac{x}{n} = \frac{68}{320} = 0.2125
	\]
	Large-sample conditions:
	\[
	n\hat{p} = 320(0.2125) = 68 \ge 10, \qquad n(1-\hat{p}) = 320(0.7875) = 252 \ge 10
	\]
	
	\textbf{Step 2: Compute the standard error of the proportion.}
	\[
	SE = \sqrt{\frac{\hat{p}(1-\hat{p})}{n}} = \sqrt{\frac{0.2125(0.7875)}{320}} \approx 0.0229
	\]
	
	\textbf{Step 3: Compute the margin of error (98\% confidence, \(z^* = 2.326\)).}
	\[
	E = z^* \cdot SE = 2.326(0.0229) \approx 0.0532
	\]
	
	\textbf{Step 4: Construct the confidence interval.}
	\[
	\hat{p} \pm E = 0.2125 \pm 0.0532 \Rightarrow [0.159,\ 0.266]
	\]
	
	\textbf{Step 5: Interpret the confidence interval in context.}
	We are 98\% confident that between about 15.9\% and 26.6\% of all export shipments fail the quality standard.
	
	\section*{5. Financial Literacy Completion}
		
	A national education board evaluates \textbf{500 students} who completed a financial literacy module. In the sample, \textbf{275 students} pass the final assessment.
	
	Construct and interpret a \textbf{99\% confidence interval} for the population proportion of students who pass the assessment.
	
	\subsubsection*{(C7) Constructs the confidence interval for a proportion}
	To construct the confidence interval for a proportion, we identify the population proportion, compute \(\hat{p}\), verify large-sample conditions, select the correct \(z^*\), compute \(SE\) and the margin of error, construct the interval, and interpret the result in context.
	
	\textbf{Step 1: Identify the population parameter and the sample proportion.}
	Let \(p\) be the population proportion of students who pass the assessment. The sample proportion is
	\[
	\hat{p} = \frac{x}{n} = \frac{275}{500} = 0.55
	\]
	Large-sample conditions:
	\[
	n\hat{p} = 500(0.55) = 275 \ge 10, \qquad n(1-\hat{p}) = 500(0.45) = 225 \ge 10
	\]
	
	\textbf{Step 2: Compute the standard error of the proportion.}
	\[
	SE = \sqrt{\frac{\hat{p}(1-\hat{p})}{n}} = \sqrt{\frac{0.55(0.45)}{500}} \approx 0.0222
	\]
	
	\textbf{Step 3: Compute the margin of error (99\% confidence, \(z^* = 2.576\)).}
	\[
	E = z^* \cdot SE = 2.576(0.0222) \approx 0.0573
	\]
	
	\textbf{Step 4: Construct the confidence interval.}
	\[
	\hat{p} \pm E = 0.55 \pm 0.0573 \Rightarrow [0.493,\ 0.607]
	\]
	
	\textbf{Step 5: Interpret the confidence interval in context.}
	We are 99\% confident that between about 49.3\% and 60.7\% of all students pass the financial literacy assessment.
	
	\section*{6. Regional Health Insurance Coverage}
		
	A health ministry combines survey results from two regions and finds that \textbf{312 of 600 adults} report having health insurance. The ministry wants to estimate the overall coverage proportion for adults in the country.
	
	Construct and interpret a \textbf{98\% confidence interval} for the population proportion of adults with health insurance.
	
	\subsubsection*{(C7) Constructs the confidence interval for a proportion}
	To construct the confidence interval for a proportion, we identify the population proportion, compute \(\hat{p}\), verify large-sample conditions, select the correct \(z^*\), compute \(SE\) and the margin of error, construct the interval, and interpret the result in context.
	
	\textbf{Step 1: Identify the population parameter and the sample proportion.}
	Let \(p\) be the population proportion of adults with health insurance. The sample proportion is
	\[
	\hat{p} = \frac{x}{n} = \frac{312}{600} = 0.52
	\]
	Large-sample conditions:
	\[
	n\hat{p} = 600(0.52) = 312 \ge 10, \qquad n(1-\hat{p}) = 600(0.48) = 288 \ge 10
	\]
	
	\textbf{Step 2: Compute the standard error of the proportion.}
	\[
	SE = \sqrt{\frac{\hat{p}(1-\hat{p})}{n}} = \sqrt{\frac{0.52(0.48)}{600}} \approx 0.0204
	\]
	
	\textbf{Step 3: Compute the margin of error (98\% confidence, \(z^* = 2.326\)).}
	\[
	E = z^* \cdot SE = 2.326(0.0204) \approx 0.0474
	\]
	
	\textbf{Step 4: Construct the confidence interval.}
	\[
	\hat{p} \pm E = 0.52 \pm 0.0474 \Rightarrow [0.473,\ 0.567]
	\]
	
	\textbf{Step 5: Interpret the confidence interval in context.}
	We are 98\% confident that between about 47.3\% and 56.7\% of all adults in the country have health insurance.
	
	\section*{7. Food Safety Inspection Compliance}
		
	A national agency inspects \textbf{850 food processing facilities} to determine compliance with a new safety standard. In the sample, \textbf{102 facilities} are found to be noncompliant.
	
	Construct and interpret a \textbf{99\% confidence interval} for the population proportion of facilities that are noncompliant with the safety standard.
	
	\subsubsection*{(C7) Constructs the confidence interval for a proportion}
	To construct the confidence interval for a proportion, we identify the population proportion, compute \(\hat{p}\), verify large-sample conditions, select the correct \(z^*\), compute \(SE\) and the margin of error, construct the interval, and interpret the result in context.
	
	\textbf{Step 1: Identify the population parameter and the sample proportion.}
	Let \(p\) be the population proportion of facilities that are noncompliant. The sample proportion is
	\[
	\hat{p} = \frac{x}{n} = \frac{102}{850} = 0.12
	\]
	Large-sample conditions:
	\[
	n\hat{p} = 850(0.12) = 102 \ge 10, \qquad n(1-\hat{p}) = 850(0.88) = 748 \ge 10
	\]
	
	\textbf{Step 2: Compute the standard error of the proportion.}
	\[
	SE = \sqrt{\frac{\hat{p}(1-\hat{p})}{n}} = \sqrt{\frac{0.12(0.88)}{850}} \approx 0.0111
	\]
	
	\textbf{Step 3: Compute the margin of error (99\% confidence, \(z^* = 2.576\)).}
	\[
	E = z^* \cdot SE = 2.576(0.0111) \approx 0.0287
	\]
	
	\textbf{Step 4: Construct the confidence interval.}
	\[
	\hat{p} \pm E = 0.12 \pm 0.0287 \Rightarrow [0.091,\ 0.149]
	\]
	
	\textbf{Step 5: Interpret the confidence interval in context.}
	We are 99\% confident that between about 9.1\% and 14.9\% of all food processing facilities are noncompliant with the safety standard.
	
\end{document}
